\section{Introduction}
\label{sec:introduction}

Game theory has long been a stable consumer of formal-mathematical
notation, but its primary objects — strategy spaces, equilibria,
mechanisms — are presented in textbooks~\citep{OsborneRubinstein1994,
MaschlerSolanZamir2013, ShohamLeytonBrown2008, Myerson1991} in a
mixture of styles that resists direct mechanization.  A
mechanization must commit, line by line, to whether ``a game'' means
a tree, a payoff matrix, an information-set diagram, a stochastic
process, or a Bayesian belief structure; whether ``a strategy'' is a
function on histories, on information sets, or on observation
sequences; whether ``probability'' lives in the discrete or the
measure-theoretic world; and whether ``equilibrium'' is a predicate
on a single profile, a profile distribution, or a correlation
device.  The textbook style routinely picks the most convenient
answer locally, then quietly switches when the next theorem
benefits from a different one.

This chapter describes a Lean~4 library that picks one answer at
the bottom and pays the cost of plumbing every other notion through
it.  At the bottom is a finite, discrete-probability, utility-aware
object — the \emph{kernel game}.  Every other representation in the
library — normal-form, extensive-form, multi-agent influence
diagrams, multi-round games, factored-observation stochastic games
— is a \emph{language} whose semantics is a kernel game.  Solution
concepts are defined once on the kernel-game (and game-form) layer
and inherited by every language.  Theorems about specific
representations are stated either as theorems about their kernel
games or as theorems about commuting compilations between languages.

\subsection{Contributions}

\begin{enumerate}[leftmargin=*]

\item \emph{A semantic core that is small enough to share.}  The
\TT{KernelGame} record (\S\ref{sec:semantic-core}) and its
utility-free shadow \TT{GameForm} are tiny:\ a strategy family, an
outcome type, a stochastic kernel, and (for kernel games) a utility.
Solution concepts are defined as predicates on this carrier and on
\TT{GameForm}, parameterized by an arbitrary preference relation on
outcome distributions.  Every representation downstream — and every
external consumer of the library — works against the same predicate
schema.

\item \emph{A protocol/preference split that pays for itself.}
Following classical mechanism-design practice~\citep{Myerson1991},
we separate the strategy-and-outcome structure (the \emph{game
form}) from the payoff function.  Lean lets us make the split
structural rather than rhetorical:\ \TT{GameForm} is a record without
\TT{utility}, and \TT{KernelGame} = \TT{GameForm} + \TT{utility} via
a definitional adjunction.  This lets us state Kuhn's theorem,
information-set machinery, and unilateral deviation as game-form
results — independent of preferences — and instantiate Nash,
dominance, correlated equilibrium, and Pareto efficiency for any
preference relation on outcome distributions.  Expected-utility
preferences come for free as a special case.

\item \emph{Languages and bridges as commuting compilations.}  Each
language layer (NFG, EFG, MAID, MultiRound, FOSG) ships with
syntax, an SOS-style operational semantics, and a denotational
\TT{toKernelGame} that is the official semantic anchor.  Bridges
between languages are not merely translations between syntactic
objects but theorems that two compilation paths to the same kernel
game agree.  This is what lets, for example, an EFG strategy theorem
be transported to a MAID by composing two bridges, without
re-proving anything strategic.  When a bridge enriches outcomes with
traces or changes the strategy representation, the transport
condition is phrased as observed-law deviation simulation:\ every
target-side deviation must be matched by a source-side deviation with
the same observed outcome law.

\item \emph{Classical theorems, mechanized.}  Existence of mixed
Nash equilibria (via Brouwer on the product simplex, lifted from
mathlib's fixed-point theory), of correlated and coarse correlated
equilibria, the minimax theorem, the one-shot deviation principle,
Zermelo's backward induction, and Kuhn's
behavioral-vs-mixed-strategy equivalence are all stated and proved
on the \TT{KernelGame} / \TT{GameForm} layer.  Their statements live
once; their consequences for any specific language layer are
corollaries.

\item \emph{An expressiveness layer.}  Beyond raw compilation, the
library exposes a \emph{relation-parameterized expressiveness
framework}:\ a comparison between two languages is a translation
on syntax paired with a soundness theorem stating that the
compiled kernel games agree under a chosen relation $R$
(EU-isomorphism, simulation, refinement, etc.;
\S\ref{sec:languages:expressiveness}).  This is what gives
substantive content to claims like ``every finite NFG is a
one-shot FOSG'' or ``every bounded MAID is an EFG up to
strategic-form refinement,'' both of which are mechanized as
$R$-reductions in this framework; comparisons that are not in the
library (e.g., MAID directly to multi-round) appear as targets in
\S\ref{sec:future}.

\item \emph{Mechanism-design and auction worked examples.}  Bayesian
games, a finite revelation-principle formulation, social-choice
machinery, and auction fragments (Vickrey, first-price, all-pay, VCG)
are formalized as consumers of the library.  The auction layer is
deliberately uneven:\ Vickrey and VCG carry dominant-strategy
truthfulness theorems, first-price records a no-dominant-strategy
result, and all-pay consists of format-specific payoff lemmas.
\end{enumerate}

\subsection{What this library is not}

It is not a measure-theoretic foundation.  All probability is
discrete, encoded by mathlib's \TT{PMF}~\citep{mathlib2020} (we
explain in \S\ref{sec:preliminaries} what is gained and lost by this
choice).  It is not a computational engine:\ many definitions are
\TT{noncomputable}, depending on classical choice for irrelevant
case splits, and we are explicit about where a constructive
implementation would diverge.  It is not a coverage of continuous
games:\ the theorem layer works under finite carrier hypotheses, and
the core expectation operator reduces to finite sums in those cases.
It is, finally, not a stand-alone proof assistant for
game-theoretic claims:\ it is a mechanized library, designed to be
imported, depended on, and extended.

A separate \emph{Definitional Supplement} accompanies this chapter.
The supplement documents the Lean signatures of every major object
and theorem with explicit textbook correspondence, so a reader
who wants to verify that ``what the library calls Nash equilibrium
is the textbook Nash equilibrium'' can do so by inspection.  The
chapter is the narrative; the supplement is the faithfulness
certificate.  Cross-references throughout this chapter point to
the relevant supplement section.

\subsection{Reading guide}

\Cref{sec:preliminaries} fixes the probabilistic and finiteness
discipline and explains why we made the design choices we made.
\Cref{sec:semantic-core} introduces \TT{GameForm}, \TT{KernelGame},
their morphisms, and the deviation-family schema that all
equilibrium notions instantiate.  \Cref{sec:solution-concepts}
catalogues the solution concepts.  \Cref{sec:languages} introduces
the language layer and its bridges.  \Cref{sec:theorems} states the
major theorems and sketches their structure.
\Cref{sec:mechanism-auctions} covers mechanism design and auctions.
\Cref{sec:discussion} collects the design tradeoffs we made
explicit.  \Cref{sec:related-work} situates the library against
prior formalization efforts and against the textbook tradition.
\Cref{sec:future} lists the natural extensions.
